I will prove $(a)$–$(c)$ completely, and then I will state precisely what happens in $(d)$.

Fix $x\in(0,1)$$, and set$
\[
y=\sqrt{1-x^2}\sin\psi,\qquad z=\sqrt{1-x^2}\cos\psi.
\]
Then
\[
x^2+y^2+z^2=1,
\]
and
\[
F(x,\psi)= -\frac{1}{4\pi}\log P_A(x,z)+C_A,
\]
where
\[
P_A(x,z)=\bigl((1-az)^2-b^2x^2\bigr)\bigl((1+az)^2-b^2x^2\bigr),
\qquad
a=\sqrt{\frac23},\quad b=\frac1{\sqrt3}.
\]

For later use, observe that on the sphere,
\[
(1-az)^2-b^2x^2
=1-2az+\frac23 z^2-\frac13x^2
=1-2az+\frac23 z^2-\frac13(1-y^2-z^2)
=(z-a)^2+\frac{y^2}{3},
\]
and similarly
\[
(1+az)^2-b^2x^2=(z+a)^2+\frac{y^2}{3}.
\]
Hence
\[
P_A(x,z)=\left((z-a)^2+\frac{y^2}{3}\right)\left((z+a)^2+\frac{y^2}{3}\right).
\]
In particular, if $\psi\in(0,\pi/2)$, then $y>0$, so
\[
P_A(x,z)>0.
\]
Thus $F$ is smooth on $(0,\pi/2)$.

---

### Proof of $(a)$

Since $\cos(-\psi)=\cos\psi$, we have
\[
z(x,-\psi)=\sqrt{1-x^2}\cos(-\psi)=\sqrt{1-x^2}\cos\psi=z(x,\psi).
\]
Therefore
\[
P_A\bigl(x,z(x,-\psi)\bigr)=P_A\bigl(x,z(x,\psi)\bigr),
\]
and hence
\[
F(x,-\psi)=F(x,\psi).
\]

So $\psi\mapsto F(x,\psi)$ is even.

---

### Proof of $(b)$

Since $\cos(\psi+\pi)=-\cos\psi$, we have
\[
z(x,\psi+\pi)=-z(x,\psi).
\]
Now
\[
P_A(x,-z)
=\bigl((1+a z)^2-b^2x^2\bigr)\bigl((1-a z)^2-b^2x^2\bigr)
=P_A(x,z),
\]
because the two factors are merely exchanged.

Therefore
\[
F(x,\psi+\pi)
=-\frac{1}{4\pi}\log P_A\bigl(x,z(x,\psi+\pi)\bigr)+C_A
=-\frac{1}{4\pi}\log P_A\bigl(x,z(x,\psi)\bigr)+C_A
=F(x,\psi).
\]

So $\psi\mapsto F(x,\psi)$ is $\pi$-periodic.

---

### Proof of $(c)$

For fixed $x$, the only $\psi$-dependence in $F(x,\psi)=G_A(x,z(x,\psi))$ comes from $z$. Hence, by the chain rule and the formula from the setup,
\[
\partial_\psi F(x,\psi)
=\partial_z G_A(x,z)\,\partial_\psi z
=\frac{2z(3+x^2-2z^2)}{9\pi P_A(x,z)}\cdot\bigl(-\sqrt{1-x^2}\sin\psi\bigr).
\]
So
\[
\partial_\psi F(x,\psi)
=
-\frac{2z\sqrt{1-x^2}\sin\psi\,(3+x^2-2z^2)}{9\pi P_A(x,z)}.
\]

Now let $\psi\in(0,\pi/2)$. Then

- $\sqrt{1-x^2}>0$ because $x\in(0,1)$,
- $\sin\psi>0$,
- $\cos\psi>0$, hence
  \[
  z=\sqrt{1-x^2}\cos\psi>0,
  \]
- $P_A(x,z)>0$ by the preliminary factorization above.

It remains to check that
\[
3+x^2-2z^2>0.
\]
Since
\[
z^2=(1-x^2)\cos^2\psi<1-x^2,
\]
we get
\[
3+x^2-2z^2
>
3+x^2-2(1-x^2)
=
1+3x^2
>0.
\]

Thus every factor in the numerator is positive except for the leading minus sign, while the denominator is positive. Therefore
\[
\partial_\psi F(x,\psi)<0
\qquad\text{for all }x\in(0,1),\ \psi\in(0,\pi/2).
\]

Hence, for each fixed $x\in(0,1)$, the map $\psi\mapsto F(x,\psi)$ is strictly decreasing on $(0,\pi/2)$.

---

## Discussion of $(d)$

Part $(d)$ is **false as stated**.

Indeed, using the formulas from the setup,
\[
\partial_x F(x,\psi)
=\partial_x G_A(x,z)+\partial_z G_A(x,z)\,\partial_x z,
\]
with
\[
\partial_x z=-\frac{x\cos\psi}{\sqrt{1-x^2}}.
\]
Therefore
\[
\partial_x F(x,\psi)
=
\frac{x(3+2z^2-x^2)}{9\pi P_A(x,z)}
+
\frac{2z(3+x^2-2z^2)}{9\pi P_A(x,z)}
\left(-\frac{x\cos\psi}{\sqrt{1-x^2}}\right).
\]
Since $z=\sqrt{1-x^2}\cos\psi$, this simplifies to
\[
\partial_x F(x,\psi)
=
\frac{x}{9\pi P_A(x,z)}
\Bigl[(3+2z^2-x^2)-2\cos^2\psi\,(3+x^2-2z^2)\Bigr].
\]
Now set
\[
c=\cos^2\psi\in(0,1),
\qquad z^2=(1-x^2)c.
\]
Then
\[
\partial_x F(x,\psi)
=
\frac{x}{9\pi P_A(x,z)}\,f_x(c),
\]
where
\[
f_x(c)=4(1-x^2)c^2-4(1+x^2)c+(3-x^2).
\]

Since $x>0$ and $P_A(x,z)>0$ on $(0,\pi/2)$, the sign of $\partial_x F(x,\psi)$ is exactly the sign of $f_x(c)$.

Now
\[
f_x(0)=3-x^2>0,
\qquad
f_x(1)=3-9x^2.
\]

So if
\[
x>\frac1{\sqrt3},
\]
then
\[
f_x(1)=3-9x^2<0.
\]
Because $f_x$ is continuous, this implies that
\[
f_x(c)<0
\]
for all $c$ sufficiently close to $1$, i.e. for all $\psi>0$ sufficiently small. Hence
\[
\partial_x F(x,\psi)<0
\]
for all sufficiently small $\psi\in(0,\pi/2)$.

Therefore the claim
\[
\partial_x F(x,\psi)>0
\qquad\text{for all }x\in(0,1),\ \psi\in(0,\pi/2)
\]
cannot be true.

In fact, one can sharpen the discussion:

- If $0<x<1/\sqrt3$, then the discriminant of $f_x$ is
  \[
  \Delta_x
  =16(1+x^2)^2-16(1-x^2)(3-x^2)
  =32(3x^2-1)<0.
  \]
  Since the leading coefficient $4(1-x^2)$ is positive, $f_x(c)>0$ for all $c\in\mathbb R$, hence
  \[
  \partial_x F(x,\psi)>0
  \qquad\text{for all }\psi\in(0,\pi/2).
  \]

- If $x=1/\sqrt3$, then
  \[
  f_x(c)=\frac83(1-c)^2,
  \]
  so for $c\in(0,1)$,
  \[
  f_x(c)>0,
  \]
  and again
  \[
  \partial_x F(x,\psi)>0
  \qquad\text{for all }\psi\in(0,\pi/2).
  \]

- If $x>1/\sqrt3$, then $\partial_x F(x,\psi)$ is negative for $\psi$ sufficiently close to $0$, while
  \[
  f_x(0)=3-x^2>0,
  \]
  so it is positive for $\psi$ sufficiently close to $\pi/2$. Thus the sign really does change.

So the correct global statement is:
\[
\partial_x F(x,\psi)>0\ \text{for all }\psi\in(0,\pi/2)
\quad\Longleftrightarrow\quad
0<x\le \frac1{\sqrt3}.
\]

---

Thus:

- $(a)$ is true,
- $(b)$ is true,
- $(c)$ is true,
- $(d)$ is false without an additional restriction; it becomes true if one assumes
  \[
  0<x\le \frac1{\sqrt3}.
  \]

